\section{Implementierung eines Umlagerungsprozesses}
%%%%%%%%%%%%%%%%%%%%%%%%%%%%%%%%%%%%%%%%%%%%%%%%%%%%%%%%%%%%%%%%%%%%%%%%%%%%%%%%%%%%%%%%%%%%%%%%%%%%%%%%%%%%%%%%%%%%%%%%
\subsection{Definition "Umlagerung" im Kontext der Intralogistik}
In der Intralogistik bezeichnet eine Umlagerung die interne Verlagerung von Mate
rial oder Ladeeinheiten innerhalb eines Lagersystems, ohne dass ein Warenein
oder-ausgang erfolgt. Sie dient nicht der Erfüllung eines Kundenauftrags, sondern
der Optimierung der Lagerstruktur und der Betriebsabläufe. Typische Gründe für
Umlagerungen sind:
• Platzoptimierung: Zusammenführen von Teilpaletten oder das Freiräumen be
stimmter Stellplätze für neue Einlagerungen.
• Reduktion der Zykluszeiten: Lagerware, die weit weg von der Entnahmestelle
eingelagert wurde, wird in Totzeiten des Systems näher an die Entnahmestelle
umgelagert. Dies verkürzt die Zykluszeit bei späterer Anfrage der Ware.
• Strategische Positionierung: Umlagerung von Artikeln in Bereiche mit höherer
Zugriffshäufigkeit, um Kommissionierzeiten zu reduzieren.
• Qualitätssicherung: Verbringen von Waren in Quarantäne- oder Prüfzonen.
• Technische Anforderungen: Freihalten von Stellplätzen für Wartungsarbeiten
oder zur Vermeidung von Blockaden.
• Bestandskorrekturen: Anpassungen nach Inventuren oder bei fehlerhaften Bu
chungen.
Im Gegensatz zu Ein- und Auslagerungen, die externe Materialflüsse abbilden, ist die
Umlagerung ein rein interner Prozess.



\subsection{Prozessablauf der Umlagerung}
In diesem Kapitel wird der Funktionsablauf aus einer systemorientierten Perspekti
ve betrachtet. Ziel ist es, die Interaktion der beteiligten Komponenten sowie die
Abhängigkeiten zwischen mechanischen Bewegungen, Steuerungslogik und
Sicherheitsmechanismen darzustellen. Darüber hinaus werden die wesentlichen
Prozessschritte erläutert, die für einen störungsfreien Betrieb erforderlich sind. Die
Darstellung erfolgt unabhängig von der konkreten Programmierung, um eine klare
Trennung zwischen konzeptioneller Beschreibung und technischer Umsetzung zu
gewährleisten.
Automatisiertes Modell- Hochregallager
34
4 Implementierung eines Umlagerungsprozesses
Entwurf der Schrittkette für Umlagerungen
Die im Rahmen dieser Arbeit implementierte Umlagerungsfunktion basiert konzeptionell
auf den bereits vorhandenen Schrittketten für die Einlagerungs- und Auslagerungspro
zesse. Beide Abläufe bilden die Grundlage für die Umlagerung, da sie die wesentlichen
Bewegungs- und Steuerungslogiken enthalten, die für den sicheren Transport von Lade
einheiten erforderlich sind. Die Umlagerung kombiniert diese beiden Prozessketten in
einer strukturierten Sequenz: Zunächst wird die Ladeeinheit am Quellplatz aufgenom
men (analog zur Auslagerung), anschließend erfolgt die Verfahrbewegung zum Zielplatz
und die Ablage (analog zur Einlagerung).
Allerdings wurde für die Realisierung der Umlagerungsfunktion eine eigenständige
Schrittkette entwickelt. Dies war notwendig, da der Umlagerungsprozess im Vergleich
zu den bestehenden Abläufen zusätzliche Anforderungen stellt: Zum einen sind mehr
Schritte erforderlich, um die vollständige Sequenz von Aufnahme, Transport und Ablage
abzubilden. Zum anderen müssen an zwei Stellen Verzweigungen integriert werden,
um unterschiedliche Prozesspfade abhängig von den Betriebsbedingungen zu ermög
lichen. Durch die Erstellung einer eigenen Schrittkette wird sichergestellt, dass die
Logik klar strukturiert bleibt und die zusätzlichen Schritte sowie Verzweigungen ohne
Beeinträchtigung der bestehenden Programme umgesetzt werden können.
Automatisiertes Modell- Hochregallager
35
4 Implementierung eines Umlagerungsprozesses
Besonderheiten der Umlagerungsschrittkette
Die Umlagerungsschrittkette stellt eine Erweiterung der bestehenden Abläufe für Ein
lagerung und Auslagerung dar, basiert jedoch nicht ausschließlich auf deren Struktur.
Während die Grundlogik der Bewegungen übernommen wurde, erfordert die Umla
gerung zusätzliche Schritte und Verzweigungen, um die komplexeren Anforderungen
dieses Prozesses abzubilden.
Abbildung 4.20: GRAFCET-Ausschnitt der Umlagerungsschrittkette mit zwei Verzwei
gungen zur Richtungsbestimmung
Im Gegensatz zu den linearen Schrittketten der Ein- und Auslagerung muss die Um
lagerung sowohl die horizontale als auch die vertikale Fahrtrichtung dynamisch be
stimmen. Dies wird im dargestellten GRAFCET deutlich: Nach der Initialisierung (Be
reit\_Zum\_Verfahren\_In\_X) erfolgt die erste Verzweigung, bei der abhängig von der
Soll-/Ist-Position die Fahrtrichtung der X-Achse gewählt wird (Motor\_X\_rechts oder
Motor\_X\_links). Nach Erreichen der Zielposition wird analog die Y-Achse angesteuert,
wobei ebenfalls eine Verzweigung zwischen Motor\_Y\_hoch und Motor\_Y\_runter erfolgt.
Diese beiden Verzweigungen sind notwendig, um die Umlagerung flexibel zwischen
beliebigen Lagerplätzen auszuführen. Aufgrund der zusätzlichen Schritte und der lo
gischen Verzweigungen wurde eine eigenständige Schrittkette entwickelt, anstatt die
bestehenden Abläufe lediglich zu kombinieren. Dadurch bleibt die Programmstruktur
übersichtlich, und die Steuerungslogik kann die erhöhten Anforderungen hinsichtlich

\subsection{Umsetzung im TIA Portal}
Im Rahmen der Automatisierung der Modellhochregalanlage wurde eine Umlagerfunk
tion implementiert, die den Transport von Ladeeinheiten zwischen unterschiedlichen
Lagerplätzen ermöglicht. Zur Realisierung dieser Funktion wurde eine Schrittkette ein
gesetzt, da sie eine strukturierte und übersichtliche Abbildung sequenzieller Abläufe
erlaubt. Die Schrittkette dient hierbei als Steuerungslogik, um die einzelnen Prozess
schritte– vom Anfahren des Quellplatzes über die Aufnahme der Ladeeinheit bis hin
zur Ablage am Zielplatz– in einer definierten Reihenfolge auszuführen. Durch diese
Vorgehensweise wird eine hohe Prozesssicherheit gewährleistet und die Komplexität
der Steuerung reduziert.
4.3.1 Aufruf und Funktionsprinzip der Positionsvergleichsfunktion
Abbildung 4.21: Aufruf des Positionsvergleichs
Zur Bestimmung der für die Umlagerung notwendigen Fahrwege wird eine Positions
vergleichsfunktion (Vergleiche\_Fahrweg) aufgerufen, welche aus den Sollwerten der
Quell- (P1) und Zielposition (P2) die Differenzen in den kartesischen Achsen berechnet
und daraus die jeweilige Fahrtrichtung ableitet. Der Aufruf erfolgt in Abhängigkeit des
aktuellen Umlagerungskontextes (manueller Modus und Übersichts-Modus), wobei die
Eingänge als Sollwertpaare für X und Y übergeben werden und die Ausgänge als
Richtungsflags bereitgestellt werden.
Automatisiertes Modell- Hochregallager
37
4 Implementierung eines Umlagerungsprozesses
Formal seien die Sollwerte der Positionen gegeben durch
P1 = (X1, Y1),
P2 = (X2, Y2).
Die Funktion berechnet die Restwege (Differenzen)

\iffalse
∆X =X2−X1, ∆Y =Y2−Y1,

und leitet hieraus die Richtungsentscheidung für die Verfahrbewegungen ab:

b\_Fahrtrichtung\_X =

TRUE falls ∆X > 0 (Fahrt nach rechts)
FALSE falls ∆X < 0 (Fahrt nach links)


b\_Fahrtrichtung\_Y =

TRUE falls ∆Y > 0 (Fahrt nach oben)
FALSE falls ∆Y < 0 (Fahrt nach unten)
Für den Grenzfall ∆X = 0 bzw. ∆Y = 0 wird keine Bewegungsanforderung in der jewei
ligen Achse generiert, womit die Schrittfolge unmittelbar zur nächsten Achsenbewegung
bzw. zum nachfolgenden Prozessschritt übergeht.
Der Funktionsaufruf ist in zwei Varianten ausgeführt, die jeweils an den aktiven Mo
dus (b_Merker_m bzw. b_Merker_ue) gebunden sind. In beiden Fällen werden als Ein
gangsparameter die Sollwerte von Quell- und Zielplatz übergeben; die Ausgänge setzen
die Richtungsflags:
Die Ausgabeparameter sind in beiden Fällen identisch und werden den nachgeschalte
ten GRAFCET-Verzweigungen zur Richtungswahl zugeführt:


b_Fahrtrichtung_X ⇒

Motor_X_rechts bei ∆X > 0
Motor_X_links bei ∆X < 0


b_Fahrtrichtung_Y ⇒

Motor_Y_hoch
bei ∆Y >0
Motor_Y_runter bei ∆Y < 0
Damit stellt die Positionsvergleichsfunktion die für die Verzweigungen in X- und Y-Richtung
notwendigen Entscheidungsgrößen bereit und gewährleistet eine konsistente, zustand
sabhängige Ansteuerung der Lineareinheiten. Die klare Trennung von Berechnung
(∆X,∆Y) und Aktorik-Auswahl (Richtungsflags) erhöht die Wartbarkeit und unterstützt
die modulare Erweiterbarkeit der Schrittkette.
Automatisiertes Modell- Hochregallager
38
4 Implementierung eines Umlagerungsprozesses
4.3.2 Schritt 10.1 im Kontext der Schrittkette
Die in Abbildung 4.20 dargestellte GRAFCET-Struktur verdeutlicht die logische Verzwei
gung innerhalb der Umlagerungsschrittkette. Nach dem Schritt Bereit_Zum_Verfahren_In_X
(Schritt 9) erfolgt die Entscheidung über die Fahrtrichtung der X-Achse. Diese Entschei
dung basiert auf dem Ergebnis der zuvor aufgerufenen Positionsvergleichsfunktion
(Vergleiche_Fahrweg), die die Differenz zwischen Quell- und Zielposition berechnet
und die Richtungsflags b_Fahrtrichtung_X und b_Fahrtrichtung_Y setzt.
Die Verzweigung in Schritt 10 ist wie folgt aufgebaut:
• Schritt 10.1– Motor_X_rechts: Wird das Flag b_Fahrtrichtung_X auf TRUE
gesetzt, erfolgt die Ansteuerung des Motors für die Bewegung in positiver X
Richtung.
• Schritt 10.2– Motor_X_links: Ist das Flag b_Fahrtrichtung_X FALSE, wird der
Motor für die Bewegung in negativer X-Richtung aktiviert.
Die Umsetzung in der SPS-Logik ist in Abbildung 4.22 dargestellt. Hier wird die Schritt
nummer über ein MOVE-Befehl gesetzt, sobald die Aktivierungsbedingungen erfüllt sind.
Die logische Verknüpfung erfolgt über ein Set-Reset-Glied (SR), das den Schrittstatus
verwaltet. Die Aktivierung von Schritt 10.1 ist somit direkt an die Fahrtrichtungsentschei
dung gekoppelt, die aus der Positionsvergleichsfunktion resultiert.
Abbildung 4.22: SPS-Implementierung für Schritt 10.1 (Motor_X_rechts)
Wie in der Abbildung ersichtlich, wird die Aktivierung der Schrittnummer 10.1 durch
die Transitionsbedingung T9_10A ausgelöst. Diese Bedingung wird erfüllt, sobald die
Fahrtrichtung in positiver X-Achse erforderlich ist, wie zuvor durch die Positionsver
gleichsfunktion bestimmt. Analog dazu wird bei aktiver Transitionsbedingung T9_10B die
Schrittnummer 10.2 aktiviert, wodurch die Bewegung in negativer X-Achse erfolgt.
Durch diese logische Struktur entsteht eine klare und funktionale Verzweigung inner
halb der Schrittkette. Sie ermöglicht eine dynamische Auswahl des Prozesspfades
Automatisiertes Modell- Hochregallager
39
4 Implementierung eines Umlagerungsprozesses
in Abhängigkeit der berechneten Richtungsflags und trägt somit zur Flexibilität und
Modularität des Steuerungsprogramms bei. Analog dazu wird mit Schritt 12.1 bzw. 12.2
die Fahrtrichtung der y-Achse festgelegt.
4.3.3 Lagerplatzmanagement im Umlagerungsprozess
Das Lagerplatzmanagement stellt sicher, dass die Umlagerung nur auf freie Zielplätze
erfolgt und die Bestandsdaten konsistent aktualisiert werden. Die Umsetzung erfolgt in
zwei logischen Schritten, die in den folgenden Abbildungen dargestellt sind.
Abbildung 4.23: Prüfung des Lagerplatzstatus und Setzen der Platzwahl
Die erste Abbildung zeigt die Prüfung des Lagerplatzstatus im Kontext der Schrittkette.
Sobald die Schrittnummer für die Umlagerung (i_Schrittketten_Nr_Umlagerung) den
Wert 2 erreicht, wird die Platzwahl aus der Übersicht aktiviert. Hierbei wird überprüft, ob
der Zielplatz im Datenbaustein DB_Lagerbestand als belegt oder frei gekennzeichnet ist.
Ist der Status False, wird die Variable b_Platzwahl_abgeschlossen auf TRUE gesetzt
und der Platz als frei markiert (b_Lagerplatz_belegt := FALSE). Andernfalls wird der
Platz als belegt gekennzeichnet und die Platzwahl bleibt unvollständig. Diese Logik
verhindert fehlerhafte Umlagerungen auf bereits belegte Positionen.
Abbildung 4.24: Aktualisierung des Lagerbestands nach erfolgreicher Umlagerung
Die zweite Abbildung zeigt die Aktualisierung des Lagerbestands nach erfolgreicher
Automatisiertes Modell- Hochregallager
40
4 Implementierung eines Umlagerungsprozesses
Umlagerung. Hierbei werden die Daten des Quellplatzes zurückgesetzt und die Ziel
position mit den neuen Informationen beschrieben. Dies umfasst die Statusvariable
(b_Status), die eindeutige Identifikationsnummer (di_UID_Teil), die Artikelnummer
sowie den Zeitstempel. Die Logik unterscheidet zwischen manueller Umlagerung
(b_manuelle_Umlagerung_aktiv) und Umlagerung aus der Übersicht (b_uebersicht_Umlagerung_aktiv
wobei in beiden Fällen die Datenbankeinträge konsistent angepasst werden. Durch
diese Vorgehensweise wird sichergestellt, dass die Bestandsdaten jederzeit aktuell sind
und die Nachverfolgbarkeit der Umlagerungen gewährleistet bleibt.

\subsubsection{Aufruf und Funktionsprinzip der Positionsvergleichsfunktion}
Abbildung 4.21: Aufruf des Positionsvergleichs
Zur Bestimmung der für die Umlagerung notwendigen Fahrwege wird eine Positions
vergleichsfunktion (Vergleiche_Fahrweg) aufgerufen, welche aus den Sollwerten der
Quell- (P1) und Zielposition (P2) die Differenzen in den kartesischen Achsen berechnet
und daraus die jeweilige Fahrtrichtung ableitet. Der Aufruf erfolgt in Abhängigkeit des
aktuellen Umlagerungskontextes (manueller Modus und Übersichts-Modus), wobei die
Eingänge als Sollwertpaare für X und Y übergeben werden und die Ausgänge als
Richtungsflags bereitgestellt werden.
Automatisiertes Modell- Hochregallager
37
4 Implementierung eines Umlagerungsprozesses
Formal seien die Sollwerte der Positionen gegeben durch
P1 = (X1, Y1),
P2 = (X2, Y2).
Die Funktion berechnet die Restwege (Differenzen)
∆X =X2−X1, ∆Y =Y2−Y1,
und leitet hieraus die Richtungsentscheidung für die Verfahrbewegungen ab:


b_Fahrtrichtung_X =

TRUE falls ∆X > 0 (Fahrt nach rechts)
FALSE falls ∆X < 0 (Fahrt nach links)


b_Fahrtrichtung_Y =

TRUE falls ∆Y > 0 (Fahrt nach oben)
FALSE falls ∆Y < 0 (Fahrt nach unten)
Für den Grenzfall ∆X = 0 bzw. ∆Y = 0 wird keine Bewegungsanforderung in der jewei
ligen Achse generiert, womit die Schrittfolge unmittelbar zur nächsten Achsenbewegung
bzw. zum nachfolgenden Prozessschritt übergeht.
Der Funktionsaufruf ist in zwei Varianten ausgeführt, die jeweils an den aktiven Mo
dus (b_Merker_m bzw. b_Merker_ue) gebunden sind. In beiden Fällen werden als Ein
gangsparameter die Sollwerte von Quell- und Zielplatz übergeben; die Ausgänge setzen
die Richtungsflags:
Die Ausgabeparameter sind in beiden Fällen identisch und werden den nachgeschalte
ten GRAFCET-Verzweigungen zur Richtungswahl zugeführt:


b_Fahrtrichtung_X ⇒

Motor_X_rechts bei ∆X > 0
Motor_X_links bei ∆X < 0


b_Fahrtrichtung_Y ⇒

Motor_Y_hoch
bei ∆Y >0
Motor_Y_runter bei ∆Y < 0
Damit stellt die Positionsvergleichsfunktion die für die Verzweigungen in X- und Y-Richtung
notwendigen Entscheidungsgrößen bereit und gewährleistet eine konsistente, zustand
sabhängige Ansteuerung der Lineareinheiten. Die klare Trennung von Berechnung
(∆X,∆Y) und Aktorik-Auswahl (Richtungsflags) erhöht die Wartbarkeit und unterstützt
die modulare Erweiterbarkeit der Schrittkette

\subsubsection{Schritt 10.1 im Kontext Schrittkette}
Die in Abbildung 4.20 dargestellte GRAFCET-Struktur verdeutlicht die logische Verzwei
gung innerhalb der Umlagerungsschrittkette. Nach dem Schritt Bereit_Zum_Verfahren_In_X
(Schritt 9) erfolgt die Entscheidung über die Fahrtrichtung der X-Achse. Diese Entschei
dung basiert auf dem Ergebnis der zuvor aufgerufenen Positionsvergleichsfunktion
(Vergleiche_Fahrweg), die die Differenz zwischen Quell- und Zielposition berechnet
und die Richtungsflags b_Fahrtrichtung_X und b_Fahrtrichtung_Y setzt.
Die Verzweigung in Schritt 10 ist wie folgt aufgebaut:
• Schritt 10.1– Motor_X_rechts: Wird das Flag b_Fahrtrichtung_X auf TRUE
gesetzt, erfolgt die Ansteuerung des Motors für die Bewegung in positiver X
Richtung.
• Schritt 10.2– Motor_X_links: Ist das Flag b_Fahrtrichtung_X FALSE, wird der
Motor für die Bewegung in negativer X-Richtung aktiviert.
Die Umsetzung in der SPS-Logik ist in Abbildung 4.22 dargestellt. Hier wird die Schritt
nummer über ein MOVE-Befehl gesetzt, sobald die Aktivierungsbedingungen erfüllt sind.
Die logische Verknüpfung erfolgt über ein Set-Reset-Glied (SR), das den Schrittstatus
verwaltet. Die Aktivierung von Schritt 10.1 ist somit direkt an die Fahrtrichtungsentschei
dung gekoppelt, die aus der Positionsvergleichsfunktion resultiert.
Abbildung 4.22: SPS-Implementierung für Schritt 10.1 (Motor_X_rechts)
Wie in der Abbildung ersichtlich, wird die Aktivierung der Schrittnummer 10.1 durch
die Transitionsbedingung T9_10A ausgelöst. Diese Bedingung wird erfüllt, sobald die
Fahrtrichtung in positiver X-Achse erforderlich ist, wie zuvor durch die Positionsver
gleichsfunktion bestimmt. Analog dazu wird bei aktiver Transitionsbedingung T9_10B die
Schrittnummer 10.2 aktiviert, wodurch die Bewegung in negativer X-Achse erfolgt.
Durch diese logische Struktur entsteht eine klare und funktionale Verzweigung inner
halb der Schrittkette. Sie ermöglicht eine dynamische Auswahl des Prozesspfades
Automatisiertes Modell- Hochregallager
39
4 Implementierung eines Umlagerungsprozesses
in Abhängigkeit der berechneten Richtungsflags und trägt somit zur Flexibilität und
Modularität des Steuerungsprogramms bei. Analog dazu wird mit Schritt 12.1 bzw. 12.2
die Fahrtrichtung der y-Achse festgelegt.


\subsubsection{Lagerplatzmanagement im Umlagerungsprozess}
Das Lagerplatzmanagement stellt sicher, dass die Umlagerung nur auf freie Zielplätze
erfolgt und die Bestandsdaten konsistent aktualisiert werden. Die Umsetzung erfolgt in
zwei logischen Schritten, die in den folgenden Abbildungen dargestellt sind.
Abbildung 4.23: Prüfung des Lagerplatzstatus und Setzen der Platzwahl
Die erste Abbildung zeigt die Prüfung des Lagerplatzstatus im Kontext der Schrittkette.
Sobald die Schrittnummer für die Umlagerung (i_Schrittketten_Nr_Umlagerung) den
Wert 2 erreicht, wird die Platzwahl aus der Übersicht aktiviert. Hierbei wird überprüft, ob
der Zielplatz im Datenbaustein DB_Lagerbestand als belegt oder frei gekennzeichnet ist.
Ist der Status False, wird die Variable b_Platzwahl_abgeschlossen auf TRUE gesetzt
und der Platz als frei markiert (b_Lagerplatz_belegt := FALSE). Andernfalls wird der
Platz als belegt gekennzeichnet und die Platzwahl bleibt unvollständig. Diese Logik
verhindert fehlerhafte Umlagerungen auf bereits belegte Positionen.
Abbildung 4.24: Aktualisierung des Lagerbestands nach erfolgreicher Umlagerung
Die zweite Abbildung zeigt die Aktualisierung des Lagerbestands nach erfolgreicher
Automatisiertes Modell- Hochregallager
40
4 Implementierung eines Umlagerungsprozesses
Umlagerung. Hierbei werden die Daten des Quellplatzes zurückgesetzt und die Ziel
position mit den neuen Informationen beschrieben. Dies umfasst die Statusvariable
(b_Status), die eindeutige Identifikationsnummer (di_UID_Teil), die Artikelnummer
sowie den Zeitstempel. Die Logik unterscheidet zwischen manueller Umlagerung
(b_manuelle_Umlagerung_aktiv) und Umlagerung aus der Übersicht (b_uebersicht_Umlagerung_aktiv
wobei in beiden Fällen die Datenbankeinträge konsistent angepasst werden. Durch
diese Vorgehensweise wird sichergestellt, dass die Bestandsdaten jederzeit aktuell sind
und die Nachverfolgbarkeit der Umlagerungen gewährleistet bleibt


\subsection{Visualisierung auf dem HMI}

Die Visualisierung des Umlagerungsprozesses wurde im bestehenden Stil des bisheri
gen HMI-Layouts umgesetzt, um eine konsistente Benutzerführung und ein einheitliches
Erscheinungsbild zu gewährleisten. Sämtliche grafischen Elemente, Farbschemata und
Bedienelemente orientieren sich an der zuvor implementierten Struktur für Ein- und
Auslagerungsprozesse. Ergänzend wurde das HMI um spezifische Funktionen für die
Umlagerung erweitert. Dazu zählen die Anzeige des aktuellen Quell- und Zielplatzes,
die Positionsmeldungen während der Verfahrbewegungen sowie die Möglichkeit zur
direkten Auswahl der Umlagerungsposition im Menü Übersicht.
Abbildung 4.25: Übersichtsdarstellung zur Auswahl des Umlagerplatzes im HMI
Oben abgebildet ist die Übersicht des Hochregallagers mit den einzelnen Lagerplätzen,
die in einer vertikalen Anordnung dargestellt sind. Die Platznummern (z.B. 11, 21,
31, 41) setzen sich aus den X- und Y-Koordinaten zusammen , um eine schnelle
Automatisiertes Modell- Hochregallager
41
4 Implementierung eines Umlagerungsprozesses
Orientierung zu gewährleisten. Im Zentrum befindet sich ein Popup-Bereich für die
Auswahl des Umlagerplatzes. Das Pop-Up wird geöffnet, sobald der Bediener der
Anlage einen Lagerplatz anklickt und ihn dadurch auswählt.
Abbildung 4.26: Popup-Fenster zur Anzeige von Produktinformationen und Eingabe des
neuen Lagerplatzes
Das Bild zeigt das Popup-Fenster, das zur Anzeige detaillierter Produktinformationen
sowie zum Start des neuen Umlagerprozesses dient. Auf der linken Seite werden
Statusinformationen des ausgewählten Lagerplatzes dargestellt, darunter der aktuelle
Status, die Artikelnummer, das Material, die eindeutige Identifikationsnummer (UID)
sowie ein Zeitstempel. Die rechte Seite des Fensters zeigt den zuvor in der Übersicht
ausgewählten Koordinaten für die neue Zielposition (X- und Y-Koordinaten) sowie eine
Schaltfläche zur Auslösung des Umlagerungsprozesses. Diese Struktur ermöglicht eine
klare Trennung zwischen Informationsanzeige und Steuerungsfunktion.
Automatisiertes Modell- Hochregallager
42
4 Implementierung eines Umlagerungsprozesses
Abbildung 4.27: Manuelles Menü des Umlagerungsprozesses mit Anzeige der aktuellen
Position
Abgebildet ist die Hauptansicht für die Durchführung der Umlagerung. Hier werden die
Quell- und Zielpositionen in separaten Eingabebereichen dargestellt: Der obere Bereich
(Von) dient zur Eingabe der aktuellen Position (X-Soll, Y-Soll), während der untere
Bereich (Nach) die Zielposition aufnimmt. Darunter befindet sich eine Schaltfläche zur
Auslösung des Umlagerungsprozesses. Rechts ist eine schematische Darstellung der
Anlage integriert, die den räumlichen Bezug zum Hochregallager verdeutlicht. Diese
Ansicht ermöglicht eine intuitive Bedienung und stellt alle relevanten Parameter für die
Umlagerung kompakt dar

\fi